\begin{enumerate}

\item De $k = EA/L$ se obtiene el área $A = kL/E$.

\item La masa es $m = \rho A L = \rho k L^{2}/E$.

\item Para el acero se repite el cálculo con sus datos.

\end{enumerate}

\noindent\textbf{Operaciones y resultados intermedios.}

\begin{enumerate}

\item Área necesaria con 7075.

\[\fitmath{y_{1}=\frac{50000000\cdot 2}{7.17\times10^{10}}\approx 0{,}0013947\,\mathrm{m^{2}}}\]

\item Masa del tirante de acero, para comparar.

\[\fitmath{y_{2}=7900\cdot \frac{50000000\cdot 2}{2\times10^{11}}\cdot 2\approx 7{,}9\,\mathrm{kg}}\]

\item Masa del tirante de 7075, m = rho A L.

\[\fitmath{y_{3}=2810\cdot \frac{50000000\cdot 2}{7.17\times10^{10}}\cdot 2\approx 7{,}838215\,\mathrm{kg}}\]

\end{enumerate}
